\documentclass[aspectratio=169]{beamer} 
\synctex=1
% 16:9 宽屏比例

\usepackage{../../Styles/mystyle-PPT}

\graphicspath{{./image/}}% 指定图片目录，后续可以直接使用图片文件名。

% beamer 主题、颜色与字体
\usetheme{metropolis}
\definecolor{mycolor}{RGB}{0, 176, 183}% 自定义青色
\definecolor{mycolor2}{RGB}{0, 191, 165}% 自定义亮绿色
% \usecolortheme[named=mycolor]{structure}% 让所有结构性元素（例如\item的序号等）都默认保持mycolor颜色

\usefonttheme{professionalfonts}% 让所有公式的字母、数字恢复为标准 LaTeX（Computer Modern）的衬线数学字体

\makeatother

% 演示文稿信息
\title{2020 年全国硕士研究生招生考试 \\ 数学一试题解析}
\subtitle{完整解答}
\author{邹文杰}
\institute{}
\date{\today}

\begin{document}

% 标题页
\begin{frame}
\titlepage
\end{frame}

% 目录
\begin{frame}{目录 CONTENTS}
\begin{columns}[T]
\column{0.48\textwidth}
\begin{block}{选择题}
1--8 小题
\end{block}
\begin{block}{填空题}
9--14 小题
\end{block}
\column{0.48\textwidth}
\begin{block}{解答题}
15--23 小题
\end{block}
\end{columns}
\end{frame}

% ============================================================
% 选择题部分
% ============================================================
\section{选择题}

\begin{frame}{选择题 (1--2)}
\begin{block}{第1题}
当 $x\to 0^+$ 时，下列无穷小量中最高阶的是
$\mathrm{A.}\int_{0}^{x}(e^{t^{2}} - 1)\mathrm{d}t \quad
\mathrm{B.}\int_{0}^{x}\ln (1 + \sqrt{t^{3}})\mathrm{d}t \quad
\mathrm{C.}\int_{0}^{\sin x}\sin t^{2}\mathrm{d}t \quad
\mathrm{D.}\int_{0}^{1 - \cos x}\sqrt{\sin^{3}t}\mathrm{d}t$.
\\
\textcolor{red}{\textbf{答案：}$\mathbf{D}$} \\
解析：利用等价无穷小代换，$\mathrm{D}$ 的阶为 $x^5$，最高.
\end{block}
\end{frame}

\begin{frame}
\begin{block}{第2题}
设函数 $f(x)$ 在 $(-1,1)$ 内有定义，且 $\lim\limits_{x\to 0}f(x)=0$，则
\begin{enumerate}[A.]
\item $\lim\limits_{x\to 0}\dfrac{f(x)}{\sqrt{|x|}}=0 \Longrightarrow  f(x)$ 在 $x=0$ 处可导;
\item $\lim\limits_{x\to 0}\dfrac{f(x)}{x^{2}}=0 \Longrightarrow  f(x)$ 在 $x=0$ 处可导;
\item $\lim\limits_{x\to 0}\dfrac{f(x)}{\sqrt{|x|}}=0$ 是 $f(x)$ 在 $x=0$ 处可导的必要条件;
\item $\lim\limits_{x\to 0}\dfrac{f(x)}{x^{2}}=0$ 是 $f(x)$ 在 $x=0$ 处可导的必要条件.
\end{enumerate}
\textcolor{red}{\textbf{答案：}$\mathbf{C}$} \\
解析：可导必连续且 $\lim \dfrac{f(x)}{x}$ 存在，故 $f(x)$ 是 $x$ 的高阶无穷小，从而 $\lim \dfrac{f(x)}{\sqrt{|x|}}=0$.
\end{block}
\end{frame}

\begin{frame}{选择题 (3--4)}
\begin{block}{第3题}
设 $f(x,y)$ 在 $(0,0)$ 可微，$f(0,0)=0$，$\boldsymbol{n}=\left.\left( \frac{\partial f}{\partial x},\frac{\partial f}{\partial y},-1 \right) \right|_{(0,0)}$，非零向量 $\bm{d}$ 与 $\bm{n}$ 垂直，则
\begin{enumerate}[A.]
\item $\lim_{(x,y)\to(0,0)}\dfrac{|\bm{n}\cdot(x,y,f(x,y))|}{\sqrt{x^2+y^2}}=0$;
\item $\lim\limits_{(x,y)\to(0,0)}\dfrac{|\bm{n}\times(x,y,f(x,y))|}{\sqrt{x^2+y^2}}=0$;
\item $\lim\limits_{(x,y)\to(0,0)}\dfrac{|\bm{d}\cdot(x,y,f(x,y))|}{\sqrt{x^2+y^2}}=0$;
\item $\lim\limits_{(x,y)\to(0,0)}\dfrac{|\bm{d}\times(x,y,f(x,y))|}{\sqrt{x^2+y^2}}=0$.
\end{enumerate}
\textcolor{red}{\textbf{答案：}$\mathbf{A}$} \\
解析：由可微性 $f(x,y)=f_x(0,0)x+f_y(0,0)y+o(\sqrt{x^2+y^2})$，代入即得.
\end{block}
\end{frame}

\begin{frame}
\begin{block}{第4题}
设 $R$ 为幂级数 $\sum\limits_{n=1}^{\infty} a_n x^n$ 的收敛半径，$r$ 是实数，则
\begin{enumerate}[A.]
\item $\sum\limits_{n=1}^{\infty} a_{2n} r^{2n}$ 发散 $\Longrightarrow  |r|\geqslant R$;
\item $\sum\limits_{n=1}^{\infty} a_{2n} r^{2n}$ 发散 $\Longrightarrow  |r|\leqslant R$.
\end{enumerate}
\textcolor{red}{\textbf{答案：}$\mathbf{A}$} \\
解析：若 $|r|<R$，则 $\sum a_{2n}r^{2n}$ 绝对收敛，故发散时必有 $|r|\geqslant R$.
\end{block}
\end{frame}

\begin{frame}{选择题 (5--6)}
\begin{block}{第5题}
若矩阵 $\bm{A}$ 经初等列变换化成 $\bm{B}$，则
\begin{enumerate}[A.]
\item $\exists P,\ PA=B$;
\item $\exists P,\ BP=A$;
\item $\exists P,\ PB=A$;
\item $Ax=0$ 与 $Bx=0$ 同解.
\end{enumerate}
\textcolor{red}{\textbf{答案：}$\mathbf{B}$} \\
解析：初等列变换等价于右乘可逆矩阵，故存在可逆矩阵 $P$ 使 $AP=B$，即 $BP^{-1}=A$，令 $P'=P^{-1}$ 得 $BP'=A$.
\end{block}
\end{frame}

\begin{frame}
\begin{block}{第6题}
已知直线 $L_1:\dfrac{x-a_2}{a_1}=\dfrac{y-b_2}{b_1}=\dfrac{z-c_2}{c_1}$ 与 $L_2:\dfrac{x-a_3}{a_2}=\dfrac{y-b_3}{b_2}=\dfrac{z-c_3}{c_2}$ 相交。设 $\alpha_i=(a_i,b_i,c_i)^T$，则
\begin{enumerate}[A.]
\item $\alpha_1$ 可由 $\alpha_2,\alpha_3$ 线性表示;
\item $\alpha_2$ 可由 $\alpha_1,\alpha_3$ 线性表示;
\item $\alpha_3$ 可由 $\alpha_1,\alpha_2$ 线性表示;
\item $\alpha_1,\alpha_2,\alpha_3$ 线性无关.
\end{enumerate}
\textcolor{red}{\textbf{答案：}$\mathbf{C}$} \\
解析：由相交得 $\alpha_3 = t\alpha_1 + (1-t)\alpha_2$.
\end{block}
\end{frame}

\begin{frame}{选择题 (7--8)}
\begin{block}{第7题}
设 $A,B,C$ 为随机事件，$P(A)=P(B)=P(C)=\frac14$，$P(AB)=0$，$P(AC)=P(BC)=\frac1{12}$，则 $A,B,C$ 中恰有一个发生的概率为

$\mathrm{A.}\,\,\frac34$;$\quad \mathrm{B.}\,\,\frac23$;$\quad\mathrm{C.}\,\,\frac12$;$\quad\mathrm{D.}\,\,\frac5{12}$.

\textcolor{red}{\textbf{答案：}$\mathbf{D}$} \\
解析：利用加法公式和容斥，$P(ABC)=0$，然后
$P(A\overline{B}\overline{C})+P(\overline{A}B\overline{C})+P(\overline{A}\overline{B}C)=\frac5{12}$.
\end{block}
\end{frame}

\begin{frame}
\begin{block}{第8题}
设 $X_1,\cdots,X_{100}$ 为来自 $P(X=0)=P(X=1)=\frac12$ 的总体，$\Phi(x)$ 为标准正态分布函数，则利用中心极限定理可得 $P\{\sum\limits_{i=1}^{100}X_i\leqslant 55\}$ 的近似值为

$\mathrm{A.}\,\,1-\Phi(1)$;$\quad \mathrm{B.}\,\,\Phi(1)$;$\quad\mathrm{C.}\,\,1-\Phi(2)$;$\quad\mathrm{D.}\,\,\Phi(2)$.

\textcolor{red}{\textbf{答案：}$\mathbf{B}$} \\
解析：$E(X)=\frac12, D(X)=\frac14$，故 $\sum X_i \sim N(50,25)$，$P(\sum X_i\leqslant 55)=P(\frac{\sum X_i-50}{5}\leqslant 1)\approx \Phi(1)$.
\end{block}
\end{frame}

% ============================================================
% 填空题部分
% ============================================================
\section{填空题}

\begin{frame}{填空题 (9--11)}
\begin{block}{第9题}
$\lim\limits_{x\to 0}\left[\dfrac{1}{e^x-1}-\dfrac{1}{\ln(1+x)}\right]=$ \underline{$\quad\textcolor{red}{-1}\quad$}
\end{block}

\begin{block}{第10题}
设 $\begin{cases}x=\sqrt{t^2+1}\\ y=\ln(t+\sqrt{t^2+1})\end{cases}$，则 $\left.\dfrac{d^2 y}{dx^2}\right|_{t=1}=$ \underline{$\quad\textcolor{red}{-\sqrt{2}}\quad$}
\end{block}

\begin{block}{第11题}
若 $f(x)$ 满足 $f''(x)+a f'(x)+f(x)=0$（$a>0$），$f(0)=m,\ f'(0)=n$，则 $\int_0^{+\infty} f(x)\,dx=$ \underline{$\quad \textcolor{red}{n+am}\quad$}
\end{block}
\end{frame}

\begin{frame}{填空题 (12--14)}
\begin{block}{第12题}
设 $f(x,y)=\int_0^{xy} e^{x^2}\mathrm{d}t$，则 $\left.\dfrac{\partial^2 f}{\partial x\partial y}\right|_{(1,1)}=$ \underline{$\quad \textcolor{red}{4e}\quad$}
\end{block}

\begin{block}{第13题}
行列式 $\begin{vmatrix}
2 & 0 & -1 & 1 \\
1 & 0 & -1 & a \\
0 & a & 1 & -1 \\
-1 & 1 & a & 0
\end{vmatrix} =$ \underline{$\quad \textcolor{red}{-4a^2+4} \quad$}
\end{block}

\begin{block}{第14题}
设 $X\sim U(-\frac{\pi}{2},\frac{\pi}{2})$，$Y=\sin X$，则 $\operatorname{Cov}(X,Y)=$ \underline{$\quad \textcolor{red}{\frac{2}{\pi}}\quad$}
\end{block}
\end{frame}

% ============================================================
% 解答题部分
% ============================================================


\section{解答题}

\begin{frame}{第15题}
\begin{block}{题目}
求二元函数 $f(x,y)=x^3+8y^3-xy$ 的极值.
\end{block}
\begin{block}{解答}
令
\begin{align*}
f_x=3x^2-y=0, \quad
f_y=24y^2-x=0,
\end{align*}
得驻点 $(0,0)$、$(\frac16,\frac1{12})$. 
判别 $A=f_{xx}=6x,\ B=f_{xy}=-1,\ C=f_{yy}=48y$. \\
在 $(0,0)$：$AC-B^2=-1<0$，非极值；在 $(\frac16,\frac1{12})$：$A=1,\ B=-1,\ C=4$，$AC-B^2=3>0$ 且 $A>0$，故为极小值点，极小值 $f(\frac16,\frac1{12})=-\frac1{216}$.
\end{block}
\end{frame}

\begin{frame}{第16题}
\begin{block}{题目}
计算 $I=\oint_L \dfrac{4x-y}{4x^2+y^2}\,dx + \dfrac{x+y}{4x^2+y^2}\,dy$，其中 $L$ 为 $x^2+y^2=2$，逆时针.
\end{block}
\begin{block}{解答}
$P=\dfrac{4x-y}{4x^2+y^2}, Q=\dfrac{x+y}{4x^2+y^2}$，计算 $\dfrac{\partial Q}{\partial x}-\dfrac{\partial P}{\partial y}=0$（除原点外）. \\
取小椭圆 $L_0: 4x^2+y^2=r^2$（逆时针），由格林公式知 $I=\oint_{L_0} P\,dx+Q\,dy$. \\
参数化 $x=\frac{r}{2}\cos\theta,\ y=r\sin\theta$，代入得 $I=\pi$.
\end{block}
\end{frame}

\begin{frame}{第17题}
\begin{block}{题目}
设数列 $\{a_n\}$ 满足 $a_1=1,\ (n+1)a_{n+1}=(n+\frac12)a_n$，证明当 $|x|<1$ 时幂级数 $\sum\limits_{n=1}^\infty a_n x^n$ 收敛，并求其和函数.
\end{block}
\begin{block}{解答}
由递推得 $\lim\limits_{n\to\infty} \left|\frac{a_{n+1}}{a_n}\right|=1$，故收敛半径 $R=1$. \\
令 $S(x)=\sum\limits_{n=1}^\infty a_n x^n$，则
\begin{align*}
S'(x)=1+\sum_{n=1}^\infty (n+1)a_{n+1}x^n = 1+xS'(x)+\frac12 S(x).
\end{align*}
\end{block}
\end{frame}

\begin{frame}
解得
\begin{align*}
S'(x)-\frac{1}{2(1-x)}S(x)=\frac{1}{1-x},
\end{align*}
通解 $S(x)=\dfrac{C}{\sqrt{1-x}}-2$，由 $S(0)=0$ 得 $C=2$，故 $S(x)=\dfrac{2}{\sqrt{1-x}}-2$.
\end{frame}

\begin{frame}{第18题}
\begin{block}{题目}
设 $\Sigma$ 为曲面 $z=\sqrt{x^2+y^2}$（$1\leqslant x^2+y^2\leqslant 4$）的下侧，$f(x)$ 连续，计算
\begin{align*}
I=\iint_\Sigma [x f(xy)+2x-y]\,dy\,dz + [y f(xy)+2y+x]\,dz\,dx + [z f(xy)+z]\,dx\,dy.
\end{align*}
\end{block}
\begin{block}{解答}
利用高斯公式补上顶面和底面，但曲面为锥面下侧，直接计算得 $I=\iint_D \sqrt{x^2+y^2}\,d\sigma$，其中 $D:1\leqslant x^2+y^2\leqslant 4$. \\
极坐标下
\begin{align*}
I=\int_0^{2\pi}d\theta\int_1^2 r^2 dr = \frac{14\pi}{3}.
\end{align*}
\end{block}
\end{frame}

\begin{frame}{第19题}
\begin{block}{题目}
设 $f(x)$ 在 $[0,2]$ 上连续可导，$f(0)=f(2)=0$，$M=\max\limits_{x\in[0,2]}|f(x)|$. 证明：
(1) 存在 $\xi\in(0,2)$ 使 $|f'(\xi)|\geqslant M$；
(2) 若对任意 $x\in(0,2)$，$|f'(x)|\leqslant M$，则 $M=0$.
\end{block}
\begin{block}{证明}
(1) 设 $|f(c)|=M$. 若 $0<c\leqslant1$，由拉格朗日中值定理 $f(c)=f'(\xi_1)c$，故 $M\leqslant |f'(\xi_1)|$；若 $1\leqslant c<2$，同理 $M\leqslant |f'(\xi_2)|$. \\
(2) 反证，若 $M>0$，则 $c\in(0,2)$. 当 $c\ne1$ 时，由中值定理可得矛盾；当 $c=1$ 时，令 $G(x)=f(x)-Mx$，单调递减且 $G(0)=G(1)=0$，推出 $M=0$，矛盾.
\end{block}
\end{frame}

\begin{frame}{第20题}
\begin{block}{题目}
$f(x_1,x_2)=x_1^2-4x_1x_2+4x_2^2$，经正交变换 $\binom{x_1}{x_2}=Q\binom{y_1}{y_2}$ 得 $g(y_1,y_2)=a y_1^2+4y_1y_2+b y_2^2$（$a\geqslant b$）.
\begin{enumerate}[(1)]
\item 求 $a,b$ 的值；

\item 求正交矩阵 $Q$.
\end{enumerate} 
\end{block}
\end{frame}

\begin{frame}
\begin{block}{解答}
$f$ 的矩阵 $A=\begin{pmatrix}1&-2\\-2&4\end{pmatrix}$，特征值为 $0,5$. \\
$g$ 的矩阵 $B=\begin{pmatrix}a&2\\2&b\end{pmatrix}$，因合同且正交变换，$A\sim B$，故
\begin{align*}
a+b=5, \quad 
ab=4,
\end{align*}
解得 $a=4,b=1$.
分别求 $A$ 和 $B$ 的单位特征向量，构造正交矩阵 $Q=\frac15\begin{pmatrix}-4&3\\3&4\end{pmatrix}$.
\end{block}
\end{frame}

\begin{frame}{第21题}
\begin{block}{题目}
设 $\bm{A}$ 为2阶矩阵，$\bm{P}=(\alpha,\bm{A}\alpha)$，$\alpha$ 非零且不是 $\bm{A}$ 的特征向量.
\begin{enumerate}[(1)]
\item 证明 $\bm{P}$ 可逆；
\item 若 $\bm{A}^2\alpha+\bm{A}\alpha-6\alpha=0$，求 $\bm{P}^{-1}\bm{A}\bm{P}$，并判断 $\bm{A}$ 是否相似于对角矩阵.
\end{enumerate}
\end{block}
\end{frame}

\begin{frame}
\begin{block}{解答}
\begin{enumerate}[(1)]
\item 反证：若 $\alpha,\ A\alpha$ 线性相关，则 $A\alpha=\lambda\alpha$，与条件矛盾. \\
\item 由条件 $A^2\alpha=6\alpha-A\alpha$，故
\begin{align*}
AP=A(\alpha,A\alpha)=(A\alpha,A^2\alpha)=(A\alpha,6\alpha-A\alpha)=P\begin{pmatrix}0&6\\1&-1\end{pmatrix},
\end{align*}
故 $P^{-1}AP=B=\begin{pmatrix}0&6\\1&-1\end{pmatrix}$. \\
$B$ 的特征值为 $2,-3$，故 $B$ 可对角化，从而 $A$ 相似于对角矩阵.
\end{enumerate}
\end{block}
\end{frame}

\begin{frame}{第22题}
\begin{block}{题目}
设 $X_1,X_2,X_3$ 独立，$X_1,X_2\sim N(0,1)$，$P(X_3=0)=P(X_3=1)=\frac12$，$Y=X_3X_1+(1-X_3)X_2$.
\begin{enumerate}[(1)]
\item 求 $(X_1,Y)$ 的分布函数；
\item 证明 $Y\sim N(0,1)$.
\end{enumerate}
\end{block}
\end{frame}

\begin{frame}
\begin{block}{解答}
\begin{enumerate}[(1)]
\item 由全概率公式：
\begin{align*}
F(x,y)=P(X_1\leqslant x,Y\leqslant y)=\frac12 P(X_1\leqslant x,X_2\leqslant y)+\frac12 P(X_1\leqslant x,X_1\leqslant y).
\end{align*}
当 $x<y$ 时，$F=\frac12\Phi(x)\Phi(y)+\frac12\Phi(x)$；当 $x\geqslant y$ 时，$F=\frac12\Phi(x)\Phi(y)+\frac12\Phi(y)$. \\
\item $P(Y\leqslant y)=\frac12P(X_2\leqslant y)+\frac12P(X_1\leqslant y)=\Phi(y)$，故 $Y\sim N(0,1)$.
\end{enumerate}
\end{block}
\end{frame}

\begin{frame}{第23题}
\begin{block}{题目}
设某种元件的使用寿命 $T$ 的分布函数为 $F(t)=1-e^{-\left(\frac{t}{\theta}\right)^m}$（$t>0$，$m>0$ 已知）.
\begin{enumerate}[(1)]
\item 求 $P(T>t)$ 与 $P(T>s+t|T>s)$；

\item 任取 $n$ 个元件测得寿命 $t_1,\cdots,t_n$，若 $m$ 已知，求 $\theta$ 的最大似然估计值.
\end{enumerate}
\end{block}
\end{frame}

\begin{frame}
\begin{block}{解答}
\begin{enumerate}[(1)]
\item $P(T>t)=e^{-(\frac{t}{\theta})^m}$，$P(T>s+t|T>s)=e^{\left(\frac{s}{\theta}\right)^m-\left(\frac{s+t}{\theta}\right)^m}$.

\item 似然函数
\begin{align*}
L(\theta)=m^n\theta^{-mn}(\prod t_i)^{m-1}\exp\{-\theta^{-m}\sum t_i^m\},
\end{align*}
对数求导得
\begin{align*}
\hat{\theta}=\left(\frac1n\sum_{i=1}^n t_i^m\right)^{\frac{1}{m}}.
\end{align*}
\end{enumerate}
\end{block}
\end{frame}

\begin{frame}
\centering
\vfill
\Huge \textcolor{mycolor}{感谢观看} \\
\vfill
\end{frame}

\end{document}